\documentclass[twocolumn, prb, amsmath,amssymb,amsfonts]{revtex4}
\usepackage{graphicx,hyperref, verbatim, float, url}
\usepackage{lipsum}
\usepackage[printwatermark]{xwatermark}
\usepackage{xcolor}
\usepackage{fancyvrb}
\usepackage{wasysym}
\hypersetup{colorlinks=true, urlcolor=blue}
\usepackage{times}
\usepackage{textcomp}
\usepackage{transparent}

\newwatermark[allpages, xpos=-200,ypos=-130]{\includegraphics[scale=1.75]{watermark.pdf}}


\begin{document}
\title{Kronos}
\author{P. D. Mullen, D. Lyons, R. S. Cumbee, P. C. Stancil}
\affiliation{Department of Physics and Astronomy, University of Georgia,
Athens, Georgia 30602\\}
\date{30 October 2015}

\begin{abstract}
Kronos is a series of python and fortran codes utilizing a comprehensive charge exchange database at its backbone to generate theoretical line ratios and spectra for many single electron capture systems.  Kronos' power can be found in that it is 1) very easy to use, 2) it allows for the user to specify a variety of input parameters of interest 3) it is quick, and 4) above all,  it contains a fully accessible and comprehensive database of multi-channel Landau-Zener single electron capture $n \ell S$-resolved cross sections for many projectile ion/ neutral target pairs at a wide variety of collision energies.  These cross sections can be used in other cascade models outside of Kronos or studied directly; in other words, for some users, Kronos may never be run and it serves the sole purpose of acting as a stand-alone charge exchange cross section database.  
\end{abstract}

\maketitle
\section{Introduction}
Kronos has been a collaborative effort of the Stancil Research Group (SRG).  It contains a series of fortran and python codes authored by different members of the group and contains charge exchange data computed by the Stueckelberg program set written by Mr. David Lyons and Dr. Phillip Stancil.  Kronos requires little instruction for its use and has been designed to be as user-friendly as possible.  Suggestions and comments on additions of features or comments on possible improvements are appreciated! As of now, Kronos is only known to be running successfully on Mac and Linux machines.  
\section{Contents}
The Kronos working directory contains several directories filled with various files and file types:\\
\\
\underline{\textbf{I. CXDatabase}}:
\begin{enumerate}
\item excitation energy files (*.nist.in)
\item n-resolved cross section files (*\_nres.cs)
\item $n \ell S$-resolved cross section files:
\begin{enumerate}
\item \textbf{Bare Ion Projectiles}
\begin{enumerate}
\item Low Energy Distribution (*\_lowe.cs)
\item Statistical Distribution (*\_stat.cs)
\end{enumerate}
\end{enumerate}
\begin{enumerate}
\item \textbf{Non-Bare Ion Projectiles}
\begin{enumerate}
\item No Distribution Function (*\_mclz.cs)
\end{enumerate}
\end{enumerate}
\item Cascade Files (radrathe.dat):
\begin{enumerate}
\item Contains All Transitions Considered:
\begin{enumerate}
\item $n \ell S$ initial $\rightarrow$ $n\ell S$ final
\item Photon Energies
\item Transition Probabilities (Einstein A Coefficients)
\end{enumerate}
\end{enumerate}
\end{enumerate}
\vspace*{0.6 cm}
\underline{\textbf{II. kronos.py}}
\begin{enumerate}
\item the main script to be run
\end{enumerate}
\underline{\textbf{III. README}}
\begin{enumerate}
\item Contains this document.
\end{enumerate}
\underline{\textbf{IV: Results}}
\begin{enumerate}
\item Empty until you run kronos.py
\item After Running kronos.py:
\begin{enumerate}
\item Line Ratios ASCII File (LineRatios.dat)
\item Line Ratios Tex File (LineRatios.tex)
\item Line Ratios PDF File (LineRatios.pdf)
\item Spectrum File (spectrum.dat)
\item Miscellaneous Spectrum Files (all others)
\end{enumerate}
\end{enumerate}
\underline{\textbf{V. Source Code}}
\begin{enumerate}
\item Contains files/codes called upon by kronos.py that are mandatory for proper running.  
\end{enumerate}
\section{Requirements for Kronos}
It is worth noting that many people will not need to run Kronos; to some, it may only be a source of charge exchange cross sections (raw data can be accessed in CXDatabase directory, see above).  Otherwise, see below:\\ \\
\noindent
\underline{\textbf{Requirements:}}\\
\\
You MUST have pylab installed on your machine.  I recommend obtaining pylab from \url{http://continuum.io/downloads}.  Download the pylab distribution for Python 2.7. You must also have gfortran compilers (see below):  
\newpage
\noindent
\textit{For Linux users}:
\begin{verbatim}
sudo apt-get install gfortran
\end{verbatim}
\textit{For Mac users}:
\\
\\We recommend installing these compilers from the following site:  \url{http://hpc.sourceforge.net}  gcc compilers are also included within this distribution.  Only GNU Fortran (GCC) 4.9.2 is known to successfully compile the code.  However, no other gfortran compilers have been tested. To install gcc (GCC) 4.9.2/GNU Fortran (GCC) 4.9.2 from the recommended website, (a) download the gcc-4.9-bin.tar.gz file to your Desktop, (b) remove the .gz extension (if not done automatically by your machine) with the gunzip command:
\begin{Verbatim}[commandchars=\\\{\}]
cd  \texttildelow/Desktop
gunzip gcc-4.9-bin.tar.gz
\end{Verbatim}
and finally, (c) install the compilers:
\begin{Verbatim}[commandchars=\\\{\}]
cd  \texttildelow/Desktop
sudo tar xvf gcc-4.9-bin.tar -C /
\end{Verbatim}
\noindent
Finally, pdflatex is required to automatically generate pdf files from generated *.tex files.  You will run into errors if you do not have pdflatex installed.  
\section{Running Kronos}
First, let's make sure that kronos.py can be run as an executable.  In working directory type:
\begin{verbatim}
chmod +x kronos.py
\end{verbatim}
Now, kronos.py should be an executable.  Now, to run Kronos, in the working directory type:
\begin{verbatim}
./kronos.py --help
\end{verbatim}
This brings up a help menu.  It details how to properly run the program, i.e. we must add tags/command-line-options (-I, -Z, -A, -E, etc.).  Although all of these options are listed in the help menu, we also give them in Table 1.  Note that some options are required while others are not.
\begin{table}[H]
\centering
\footnotesize
\begin{tabular}{|c|c|}
\hline 
Term & Meaning \\ 
\hline 
./kronos.py & call the script (required) \\ 
\hline
-I & specify projectile ion (required) \\ 
\hline 
-Z & specify charge of projectile ion (required) \\ 
\hline 
-A & specify target (required) \\ 
\hline
-E & specify collision energy in \textbf{eV/u} (required) \\ 
\hline
-D &  lowe or stat distribution (required for MCLZ bare ion collisions) \\
\hline 
-R & specify FWHM resolution in eV (optional) \\
\hline
-P & Specify number of data points (optional) \\ 
\hline 
-L & specify your initials (optional) \\ 
\hline
-N & specify your name (optional) \\ 
\hline 
-X & specify rmax (optional: advanced users only) \\
\hline
-M & specify CX method (optional: MCLZ, CTMC, MOCC, etc.) \\
\hline
-G & Specify yes or no to graph (optional) \\
\hline
\end{tabular} 
\end{table}
Here, we make several comments about Table 1.  1) As of now, only MCLZ data as computed by the Stueckelberg program set are contained within the CXDatabase directory.  If one adds their own data (QMOCC, AOCC, CTMC, MCLZ, etc.) it must be placed in the proper directory within the CXDatabase with a proper file name (for example, for C$^{5+}$ collisions with H via QMOCC: c5+h\_sec\_qmocc.cs).  2) If one does not specify a method with the -M flag, Kronos has a built in hierarchy for the data used in computation.  This hierarchy is given as (in order of greatest to  least preference):
\\
\\
QMOCC $\rightarrow$ MOCC $\rightarrow$ AOCC $\rightarrow$ CTMC $\rightarrow$ MCLZ \\ \\
Again, at this time, only MCLZ data is given within the CXDatabase.  3) The -D flag is required for bare ion collisions.  Specifying the option lowe applies the low energy distribution to bare ion collision cross section data.  Specifying the option stat applies the statistical distribution to bare ion collision cross section data.  \\
\\
An example of running kronos.py for C$^{5+}$ + H$_2$O CX collisions at 10 eV/u mimicking a detector FWHM resolution of 10 eV:
\begin{verbatim}
./kronos.py -I C -Z 5 -A H2O -E 10 -R 10
\end{verbatim}
Here we see that we have specified all required tags but have left out several of the optional ones -- see Table 1.  
\\
\\
An example of running kronos.py for the bare ion CX collision C$^{6+}$ + H at 1 keV/u and applying the low energy distribution to charge exchange cross section data:
\begin{verbatim}
./kronos.py -I C -Z 6 -A H -E 1000 -D lowe
\end{verbatim}
After running the program, one can view the results in the Results directory.
\section{Previous Versions}
If you are reading this document, you have downloaded the most recent version of Kronos.  
You can download the previous version of Kronos from the following Dropbox link:
\url{https://www.dropbox.com/s/muq7u20zv4youzy/Kronos_v1.zip?dl=0}
\end{document}